\section{Cyclic points producing character ambiguity}\label{sec:ambiguity}
We construct the comparisons that show the ramification bounds are sharp. The construction has two parts: an invariant Eisenstein equation produces a rational cyclic point, and an extension of rank-one representations reproduces the resulting finite torsion module.

\begin{lemma}\label{lem:extension-lift}
Let $K$ have characteristic zero, and let $\psi:G_K\to\Z_p^\times$ be finite of order prime to $p$. Every continuous extension of free $\Z/p^n$-modules
\[
 0\longrightarrow\mu_{p^n}\longrightarrow V_n
       \longrightarrow(\Z/p^n)(\psi^{-1})\longrightarrow0
\]
lifts to a height-two Barsotti--Tate group over $K$ with graded groups $\mu_{p^\infty}$ and $(\Q_p/\Z_p)(\psi^{-1})$.
\end{lemma}
\begin{proof}
Set $T=\Z_p(1)\otimes\psi$. An underlying module splitting identifies the extension with a class in $H^1(K,T/p^nT)$. Over the finite extension $K'/K$ killing $\psi$, Kummer theory identifies the untwisted group at level $r$ with $K'^\times/(K'^\times)^{p^r}$. Its transition maps are surjective. Inflation--restriction identifies the original cohomology with the invariants under $\Delta=\Gal(K'/K)$ of the suitably twisted Kummer group. Averaging by $|\Delta|^{-1}\sum_{\delta\in\Delta}\delta$ preserves surjectivity. Thus the given class lifts compatibly to every higher level.

Choose continuous cocycles for these classes. At each step the reduction differs from the previous cocycle by a coboundary; lift the correcting zero-cochain and subtract its coboundary. The resulting cocycles are exactly compatible. Their inverse limit is a continuous $T$-valued cocycle, giving an extension of free $\Z_p$-modules with the prescribed reduction. The corresponding system of finite \'etale group schemes over $K$ is a Barsotti--Tate group. Its two graded groups acquire good reduction after killing $\psi$.
\end{proof}

\begin{lemma}\label{lem:point-ambiguity}
Suppose that $K$ has algebraically closed residue field, $p\ge5$, and $G/K$ is height two with supersingular good reduction over a tame extension. Suppose that $G$ has a rational point of exact order $p^n$ and special eigenvalues $a,a^p$ on a tame generator, where $a\in\mu_{p^2-1}$ and $a\ne1$. There exist a potentially supersingular group $M$ and a group $H$ filtered by two potentially good height-one groups, with
\[
                         M[p^n]\simeq H[p^n],
\]
whose finite special characters differ.
\end{lemma}
\begin{proof}
The determinant of $T_pG$ is $\chi_p\psi$, with $\psi$ a finite character valued in $\mu_{p-1}$, and $\psi(\tau)=a^{p+1}$. Indeed its restriction to a good field is crystalline of weight one. After twisting by $\chi_p^{-1}$ it has weight zero; the rank-one case of \citet[Corollary~2.2.6, p.~26 of the author version]{Kisin2006} identifies it with the rational Tate module of a height-one \'etale group. Such a group is constant over an algebraically closed residue field. The determinant is therefore cyclotomic on the good field, and its remaining finite value is the determinant of the special isocrystal action. The same argument shows that a potentially good height-one character is finite or cyclotomic times finite.

Let $C$ be the constant cyclic group generated by the point and put $G'=G/C$. Closure of $C$ in a good model is finite flat, so $G'$ still acquires supersingular good reduction. The multiplication-by-$p^n$ exact sequence gives
\[
 0\longrightarrow C\longrightarrow G[p^n]
       \longrightarrow G'[p^n]\longrightarrow C\longrightarrow0.
\]
The first $C$ is primitive as a free cyclic submodule; its quotient in $G[p^n]$ therefore has character $\chi_p\psi$. Twisting $G'$ by $\psi^{-1}$ gives $M$ and an extension of the form in Lemma~\ref{lem:extension-lift}. Apply that lemma to obtain $H$ with the same truncation. Isogeny preserves the special rational character. The two eigenvalue multisets are
\[
               \{a^{-1},a^{-p}\},\qquad\{1,a^{-(p+1)}\}.
\]
Putting $\lambda=a^{-1}$, the difference of their traces is
\[
                 \lambda+\lambda^p-1-\lambda^{p+1}
                              =-(\lambda-1)(\lambda^p-1)\ne0.
\]
This proves the claim on the finite tame inertia quotient.
\end{proof}

\subsection{Every index above the prime-to-$p$ threshold}
\begin{theorem}\label{thm:general-arc}
Let $p\ge5$, $n\ge2$, and let $e$ be a positive multiple of $p-1$ satisfying $e>p^{2n-3}$. There is a complete field $K$ with residue field $\overline{\F}_p$, absolute index $e$, and $\mu_p\subset K$, and groups $M,H$ as in Lemma~\ref{lem:point-ambiguity} whose $p^n$-torsion is isomorphic but whose finite special characters differ. In particular $\nuMin(p,e,h)\ge n$ for every $h\ge2$.
\end{theorem}
\begin{proof}
Set $q=p^2$, $N=q-1$, $Q=q^{n-1}$ and $x_0=e/Q$. Choose an integer
\[
                          x_0<A<px_0,\qquad N\nmid A.
\]
Such an integer always exists. If $x_0<1$, take $A=1$, using $px_0>1$. If $x_0\ge1$, take $\lfloor x_0\rfloor+1$ and increase it by one when it is divisible by $N$. The resulting integer is at most $x_0+2<px_0$.

Over $W=W(\overline{\F}_p)$ use the integral universal formal group $F_u$ of Section~\ref{sec:cyclic}. Its logarithm satisfies
\begin{equation}\label{eq:homogeneous-logarithm}
 f_u(T)=\sum_{j\ge0}l_j(u)T^{p^j},\quad l_0=1,
 \qquad pl_j=l_{j-1}u^{p^{j-1}}+l_{j-2},\quad l_{-1}=0.
\end{equation}
Put $d=Ne$, $b=d-(p-1)AQ$. The chosen interval gives
\begin{equation}\label{eq:general-arc-weights}
 b>(q-p)Aq^{n-2}>0,\qquad AQ<d<NAQ,
                 \qquad b\equiv-(p-1)A\pmod N.
\end{equation}
For example the first difference is exactly $N(e-Ap^{2n-3})$.
Define
\[
 Y_0=X^A,\qquad Y_{j+1}=[p]_{X^b}(Y_j),\qquad
 P(X)=\Phi_{X^b}(Y_{n-1}),\quad \Phi_u(T)=[p]_u(T)/T.
\]
These are compositions of integral power series in $W[[X]]$.

Modulo $p$, the multiplication series begins with $uT^p+T^q$; intervening coefficients are multiples of $u$, and all later coefficients are integral. The first inequality in~\eqref{eq:general-arc-weights} makes the $T^q$ term strictly earliest at every one of the first $n-1$ compositions. Hence
\[
                        Y_j\equiv X^{Aq^j}+O(X^{Aq^j+1})\pmod p.
\]
In the final series $P$, the $uY_{n-1}^{p-1}$ term has degree $d$. All intervening terms are later, and $Y_{n-1}^N$ has degree $NAQ>d$. Thus
\[
                         P(0)=p,\qquad P\equiv X^d+O(X^{d+1})\pmod p.
\]
Weierstrass preparation makes $R=W[[X]]/(P)$ an Eisenstein discrete valuation ring of absolute index $d$, with uniformizer $x$. The inequality $AQ<d$ prevents coefficients divisible by $p$ from changing the valuations $v_R(Y_j(x))=Aq^j$. Therefore $x^A$ has exact order $p^n$ in $F_{x^b}$.

For a primitive $N$th root $\zeta\in W$, set $a=\zeta^{-A}$. The recursion~\eqref{eq:homogeneous-logarithm} proves
\[
 l_j(a^{p-1}u)=a^{p^j-1}l_j(u),\qquad
 [p]_{a^{p-1}u}(a^{-1}T)=a^{-1}[p]_u(T),\qquad
 \Phi_{a^{p-1}u}(a^{-1}T)=\Phi_u(T).
\]
For the $l_{j-2}$ summand the exponents differ by $p^j-p^{j-2}$, which is divisible by $N$. Thus these are identities of the entire power series. Since $b\equiv-(p-1)A\pmod N$, substitution $X\mapsto\zeta X$ sends $u$ to $a^{p-1}u$, every $Y_j$ to $a^{-1}Y_j$, and $P$ to itself. It induces an automorphism $\sigma(x)=\zeta x$ of $R$, of order $N$.

Let $L=\operatorname{Frac}(R)$ and $K=L^{\langle\sigma\rangle}$. The extension is tame of degree $N$, and $e_K=e$. The integral isomorphism $\sigma^*F_{x^b}\to F_{x^b}$ given by $T\mapsto aT$ satisfies the cocycle identity and fixes the marked point, because $a\sigma(x^A)=x^A$. Descent at every finite generic level produces $G/K$ with that rational point. Its special eigenvalues are $a,a^p$; they are the two embeddings of the quadratic Teichm\"uller unit acting on the Honda isocrystal.

Finally $K$ contains $\mu_p$. Write $-p=\pi_K^e u$ with $u$ a unit. Hensel's lemma supplies a $(p-1)$st root of $u$, and $p-1\mid e$, so $K$ contains a $(p-1)$st root of $-p$. Over $W[1/p]$ this generates the unique tame extension of degree $p-1$, namely $W[1/p](\mu_p)$. Since $N\nmid A$, we have $a\ne1$, and Lemma~\ref{lem:point-ambiguity} applies. Adding the same constant height-$h-2$ group proves the last assertion.
\end{proof}

\begin{corollary}\label{cor:sharp-mixed-prime}
For $p\ge5$, $n\ge2$, and $(p-1)\mid e$ with $p\nmid e$, a supersingular/height-one filtered comparison with isomorphic $p^n$-torsion and different special characters exists over some field of index $e$ if and only if $e>p^{2n-3}$.
\end{corollary}
\begin{proof}
Theorem~\ref{thm:general-arc} proves existence. Proposition~\ref{prop:mixed-obstruction} proves necessity with $v_p(e)=0$.
\end{proof}

\subsection{A third-level construction at the next wild index}
\begin{theorem}\label{thm:two-step-arc}
Let $p\ge5$ and let $e$ be a positive multiple of $p(p-1)$ with $e>p^2$. There are groups $M,H$ over a field with residue field $\overline{\F}_p$, absolute index $e$, and $\mu_p\subset K$, with $M[p^3]\simeq H[p^3]$ and different finite special characters. The group $M$ acquires supersingular good reduction, and $H$ is filtered by two potentially good height-one groups.
\end{theorem}
\begin{proof}
Choose
\[
 \frac{e}{p^3}<A<\frac{e}{p^2},\qquad N\nmid A,
 \qquad d=eN,\quad b=\frac{eN}{p}-p^2(p-1)A,
 \quad J=b+p^3A.
\]
The same integer-choice argument as in the previous theorem applies, because $e/p^2>1$. Direct subtraction gives
\begin{align*}
 b-(p^2-p)A&=N(e-p^2A)/p>0,\\
 (p^2-p)p^2A-b&=N(p^3A-e)/p>0,\\
 d&=b+(p-1)J,\qquad A<p^2A<J<d,\\
 b&\equiv-(p-1)A\pmod N.
\end{align*}
Use $Y_0=X^A$, $Y_1=[p]_{X^b}(Y_0)$, $Y_2=[p]_{X^b}(Y_1)$, and $P=\Phi_{X^b}(Y_2)$. The first displayed inequality selects the degree-$p^2$ term at the first step; the second selects the degree-$p$ term at the next. Therefore
\[
               Y_1\equiv X^{p^2A}+O(X^{p^2A+1}),\qquad
               Y_2\equiv X^J+O(X^{J+1})\pmod p.
\]
The final selected term has degree $b+(p-1)J=d$. Its degree-$p^2$ competitor is later because $b<(p^2-p)J$. Intermediate terms contain $u$ modulo $p$ and are later than the selected terms; terms beyond degree $p^2$ are later by integrality. Hence $P(0)=p$ and $P\equiv X^d+O(X^{d+1})\pmod p$.

The Eisenstein ring $R=W[[X]]/(P)$ has $v_R(p)=d$. Since $J<d$, its marked point $x^A$ has successive valuations $A,p^2A,J$ and exact order $p^3$. The congruence for $b$ permits precisely the descent $x\mapsto\zeta x$, $T\mapsto\zeta^{-A}T$ used above. Its fixed field has absolute index $e$, contains $\mu_p$, and the descended group has special eigenvalues $a,a^p$ with $a=\zeta^{-A}\ne1$. Apply Lemma~\ref{lem:point-ambiguity}.
\end{proof}

\begin{corollary}\label{cor:sharp-mixed-wild}
Among positive multiples $e$ of $p(p-1)$ with $v_p(e)=1$, a comparison of the two classes in Theorem~\ref{thm:two-step-arc} exists with isomorphic $p^3$-torsion and different special characters if and only if $e>p^2$.
\end{corollary}
\begin{proof}
Existence was just proved. Proposition~\ref{prop:mixed-obstruction}, with $n=3$ and $v_p(e)=1$, gives the necessary strict inequality.
\end{proof}

\begin{theorem}\label{thm:second-wild-exact}
For $p\ge5$, in the height-two filtration setting over complete fields with perfect residue field and $\mu_p\subset K$,
\[
                         \nuMin(p,2p(p-1),2)=3.
\]
In particular $\nuMin(5,40,2)=3$.
\end{theorem}
\begin{proof}
Theorem~\ref{thm:two-step-arc} gives failure at $p^3$-torsion. At $e=2p(p-1)$ one may take $A=1$, $b=(p-1)(p^2-2)$ and $J=b+p^3$.

For the upper bound, take a $p^4$-torsion comparison and pass to the completed maximal unramified base. Each determinant has the form $\chi_p^d\epsilon$, with $d\in\{0,1,2\}$ and finite $\epsilon\in\mu_{p-1}$. The wild order of $\chi_p$ modulo $p^4$ is at least $p^2$: the intersection of $K$ with the cyclotomic tower has wild degree at most $p$, since $v_p(e)=1$. The ratio of the two determinants can thus be trivial only when their $d$ agree; their finite factors then agree as well.

For $d=0$, the graded pieces are potentially \'etale, so their special character is the Brauer lift of the semisimplified residual character. Cartier duality gives the case $d=2$. If $d=1$, a height-two potentially good piece is either supersingular or has its connected--\'etale filtration, which descends by uniqueness. Thus the possibilities are a whole supersingular group or a filtration by two height-one groups. Two filtered groups are compared residually, since $\chi_p\equiv1\pmod p$. A mixed comparison would require $e>p^{2\cdot4-3-1}=p^4$ by Proposition~\ref{prop:mixed-obstruction}, which is false.

For the remaining pair of supersingular groups, apply Theorem~\ref{thm:common-determinant} with $a=2p(p-1)(p^2-1)$. The recurrence gives
\[
 q_1=2p-1,\qquad q_2=2p^2-1,\qquad
 q_3=2p^3-1>2p(p^2-1)=a/(p-1).
\]
The $p^4$-torsion therefore determines the special character in every case. Since $p^3$-torsion can fail, the least index is three.
\end{proof}
